%% ECON 282E -- Session 3, Deck A1: What Does It Mean to Solve a Model?
%% Budget: 16 frames (~22 minutes). Provenance: slides/SOURCES.md.
%% Every number on this deck comes from tools/figures/s03_numbers.json, key
%% "floating_point" and "vfi_base" (generated by tools/figures/s03_generate.py
%% on a compute node, 2026-09-17).

%% ECON 282E -- Foundations of Macroeconomics (UC Riverside, Fall 2026)
%% Shared Beamer preamble for every deck in slides/S01 ... slides/S10.
%%
%% Usage (a deck lives in slides/S0X/ and is compiled from that folder):
%%   %% ECON 282E -- Foundations of Macroeconomics (UC Riverside, Fall 2026)
%% Shared Beamer preamble for every deck in slides/S01 ... slides/S10.
%%
%% Usage (a deck lives in slides/S0X/ and is compiled from that folder):
%%   %% ECON 282E -- Foundations of Macroeconomics (UC Riverside, Fall 2026)
%% Shared Beamer preamble for every deck in slides/S01 ... slides/S10.
%%
%% Usage (a deck lives in slides/S0X/ and is compiled from that folder):
%%   \input{../preamble/course_preamble.tex}
%%   \decktitle{1}{A1}{Who I Am \& Why This Course}{September 24, 2026}
%%   \begin{document}
%%   \makedecktitle
%%   ...
%%
%% Adapted from Courses/AI-ECON-2026/source/shared_preamble.tex (Metropolis theme,
%% concept/caution/example/takeaway boxes, \sectiondivider). Changes for this course:
%%   * course title block (\decktitle / \makedecktitle) instead of \makebeamertitle
%%   * \zh{...} typesets Chinese (book titles) under pdflatex via CJKutf8
%%   * researchbox for the "Research in the wild" frames
%%   * section pages off; TikZ libraries and the course map loaded here
%% Build with pdflatex only (two passes); tools/build_slides.py does this for you.

\documentclass[10pt,english,aspectratio=169]{beamer}

\usepackage[T1]{fontenc}
\usepackage{textcomp}
\usepackage[utf8]{inputenc}
\usepackage{babel}
\usepackage{CJKutf8}

\usepackage{amsmath, amssymb, amsthm, graphicx}
\usepackage{booktabs, multirow, tabularx}
\usepackage{tikz}
\usetikzlibrary{positioning,arrows.meta,shapes,calc,fit,backgrounds}
\usepackage{pgfplots}
\pgfplotsset{compat=1.16}
\usepackage[most]{tcolorbox}
\tcbuselibrary{skins, breakable}
\usepackage{listings}
\usepackage{xcolor}

\lstset{
  basicstyle=\ttfamily\small,
  keywordstyle=\color{blue!80!black}\bfseries,
  commentstyle=\color{green!50!black}\itshape,
  stringstyle=\color{orange!80!black},
  backgroundcolor=\color{gray!8},
  frame=single,
  rulecolor=\color{gray!40},
  breaklines=true,
  showstringspaces=false,
  numbers=none,
}

\usetheme{metropolis}
\usefonttheme{professionalfonts}
\metroset{sectionpage=none}

\definecolor{LightGray}{RGB}{230,230,230}
\definecolor{RedTitle}{RGB}{0,0,200}
\definecolor{VeryLightGray}{RGB}{245,245,245}
\definecolor{DarkText}{RGB}{0,0,0}
\definecolor{UCRBlue}{RGB}{0,45,106}
\definecolor{UCRGold}{RGB}{241,171,0}

% Stage colours of the course arc (used by the course map and elsewhere)
\colorlet{StageTrunk}{blue!12}
\colorlet{StageBranch}{cyan!14}
\colorlet{StageMachine}{gray!18}
\colorlet{StageWall}{red!14}
\colorlet{StageTools}{orange!20}
\colorlet{StageData}{green!16}
\colorlet{StageAgents}{violet!14}

\hypersetup{bookmarks=false,breaklinks=true,pdfborder={0 0 0},colorlinks=true,
  linkcolor=DarkText,urlcolor=RedTitle}

\setbeamercolor{background canvas}{bg=white}
\setbeamercolor{title}{fg=RedTitle}
\setbeamercolor{frametitle}{bg=VeryLightGray,fg=DarkText}
\setbeamercolor{normal text}{fg=DarkText}
\setbeamercolor{alerted text}{fg=RedTitle}
\setbeamersize{text margin left=5mm, text margin right=5mm}
\addtobeamertemplate{frametitle}{}{\vspace{-0.5em}}

\setbeamertemplate{navigation symbols}{}
\setbeamertemplate{footline}{%
  \hspace{5mm}{\usebeamerfont{page number in head/foot}\color{black!45}%
  ECON 282E \,$\cdot$\, \insertshortdeck}\hfill%
  \usebeamercolor[fg]{page number in head/foot}%
  \usebeamerfont{page number in head/foot}%
  \insertframenumber\,/\,\inserttotalframenumber\kern1em\vskip2pt%
}

% ---------------------------------------------------------------------------
% Chinese text under pdflatex (book titles etc.)
\newcommand{\zh}[1]{\begin{CJK*}{UTF8}{gbsn}#1\end{CJK*}}

% ---------------------------------------------------------------------------
% Deck title block
%   \decktitle{session}{deck id}{title}{date}
\newcommand{\insertshortdeck}{}
\newcommand{\decksession}{}
\newcommand{\deckid}{}
\newcommand{\decktitletext}{}
\newcommand{\deckdate}{}
\newcommand{\decktitle}[4]{%
  \renewcommand{\decksession}{#1}%
  \renewcommand{\deckid}{#2}%
  \renewcommand{\decktitletext}{#3}%
  \renewcommand{\deckdate}{#4}%
  \renewcommand{\insertshortdeck}{Session #1 \,$\cdot$\, Deck #1.#2}%
  \title{#3}%
  \author{Zhigang Feng}%
  \date{#4}%
}
\newcommand{\makedecktitle}{%
  \begin{frame}[plain,noframenumbering]
    \begin{tikzpicture}[remember picture,overlay]
      \fill[UCRBlue] (current page.north west) rectangle ([yshift=-0.55cm]current page.north east);
      \fill[UCRGold] ([yshift=-0.55cm]current page.north west) rectangle ([yshift=-0.62cm]current page.north east);
      \node[anchor=west, text=white, font=\small\bfseries] at ([xshift=5mm,yshift=-0.275cm]current page.north west)
        {ECON 282E \,$\cdot$\, Foundations of Macroeconomics};
      \node[anchor=east, text=white, font=\small] at ([xshift=-5mm,yshift=-0.275cm]current page.north east)
        {UC Riverside \,$\cdot$\, Fall 2026};
    \end{tikzpicture}
    \vspace*{1.1cm}
    {\usebeamerfont{title}\color{black!55}\large Session \decksession\ \,$\cdot$\, Deck \decksession.\deckid\par}
    \vspace{0.25cm}
    {\usebeamerfont{title}\color{RedTitle}\LARGE\bfseries \decktitletext\par}
    \vspace{0.7cm}
    {\large Zhigang Feng\par}
    \vspace{0.15cm}
    {\color{black!65}\deckdate\par}
    \vfill
    {\footnotesize\itshape\color{black!55}Build it on paper $\;\to\;$ solve it on the machine $\;\to\;$ push it to the frontier with AI\par}
  \end{frame}}

% ---------------------------------------------------------------------------
% Section divider (plain frame, not counted)
\newcommand{\sectiondivider}[2]{%
\begin{frame}[plain,noframenumbering]
\begin{center}
\vfill
{\Large\textbf{#1}\par}
\vspace{0.35cm}
{\normalsize #2\par}
\vfill
\end{center}
\end{frame}
}

% ---------------------------------------------------------------------------
% Boxes
\newtcolorbox{conceptbox}[1][]{colback=blue!3, colframe=RedTitle!55, boxrule=0.35mm,
  arc=1mm, left=2mm, right=2mm, top=1mm, bottom=1mm, #1}
\newtcolorbox{cautionbox}[1][]{colback=red!3, colframe=red!50!black, boxrule=0.35mm,
  arc=1mm, left=2mm, right=2mm, top=1mm, bottom=1mm, #1}
\newtcolorbox{examplebox}[1][]{colback=green!3, colframe=green!45!black, boxrule=0.35mm,
  arc=1mm, left=2mm, right=2mm, top=1mm, bottom=1mm, #1}
\newtcolorbox{takeawaybox}[1][]{colback=VeryLightGray, colframe=RedTitle!55, boxrule=0.35mm,
  arc=1mm, left=2mm, right=2mm, top=1mm, bottom=1mm, #1}
\newtcolorbox{researchbox}[1][]{enhanced, colback=violet!4, colframe=violet!60!black,
  boxrule=0.35mm, arc=1mm, left=2mm, right=2mm, top=1mm, bottom=1mm,
  fonttitle=\bfseries\small, title={Research in the wild}, #1}

\newcommand{\compacttables}{\renewcommand{\arraystretch}{1.15}}
\newcommand{\src}[1]{{\tiny\color{black!55}#1}}

% ---------------------------------------------------------------------------
\input{../preamble/course_map.tex}

%%   \decktitle{1}{A1}{Who I Am \& Why This Course}{September 24, 2026}
%%   \begin{document}
%%   \makedecktitle
%%   ...
%%
%% Adapted from Courses/AI-ECON-2026/source/shared_preamble.tex (Metropolis theme,
%% concept/caution/example/takeaway boxes, \sectiondivider). Changes for this course:
%%   * course title block (\decktitle / \makedecktitle) instead of \makebeamertitle
%%   * \zh{...} typesets Chinese (book titles) under pdflatex via CJKutf8
%%   * researchbox for the "Research in the wild" frames
%%   * section pages off; TikZ libraries and the course map loaded here
%% Build with pdflatex only (two passes); tools/build_slides.py does this for you.

\documentclass[10pt,english,aspectratio=169]{beamer}

\usepackage[T1]{fontenc}
\usepackage{textcomp}
\usepackage[utf8]{inputenc}
\usepackage{babel}
\usepackage{CJKutf8}

\usepackage{amsmath, amssymb, amsthm, graphicx}
\usepackage{booktabs, multirow, tabularx}
\usepackage{tikz}
\usetikzlibrary{positioning,arrows.meta,shapes,calc,fit,backgrounds}
\usepackage{pgfplots}
\pgfplotsset{compat=1.16}
\usepackage[most]{tcolorbox}
\tcbuselibrary{skins, breakable}
\usepackage{listings}
\usepackage{xcolor}

\lstset{
  basicstyle=\ttfamily\small,
  keywordstyle=\color{blue!80!black}\bfseries,
  commentstyle=\color{green!50!black}\itshape,
  stringstyle=\color{orange!80!black},
  backgroundcolor=\color{gray!8},
  frame=single,
  rulecolor=\color{gray!40},
  breaklines=true,
  showstringspaces=false,
  numbers=none,
}

\usetheme{metropolis}
\usefonttheme{professionalfonts}
\metroset{sectionpage=none}

\definecolor{LightGray}{RGB}{230,230,230}
\definecolor{RedTitle}{RGB}{0,0,200}
\definecolor{VeryLightGray}{RGB}{245,245,245}
\definecolor{DarkText}{RGB}{0,0,0}
\definecolor{UCRBlue}{RGB}{0,45,106}
\definecolor{UCRGold}{RGB}{241,171,0}

% Stage colours of the course arc (used by the course map and elsewhere)
\colorlet{StageTrunk}{blue!12}
\colorlet{StageBranch}{cyan!14}
\colorlet{StageMachine}{gray!18}
\colorlet{StageWall}{red!14}
\colorlet{StageTools}{orange!20}
\colorlet{StageData}{green!16}
\colorlet{StageAgents}{violet!14}

\hypersetup{bookmarks=false,breaklinks=true,pdfborder={0 0 0},colorlinks=true,
  linkcolor=DarkText,urlcolor=RedTitle}

\setbeamercolor{background canvas}{bg=white}
\setbeamercolor{title}{fg=RedTitle}
\setbeamercolor{frametitle}{bg=VeryLightGray,fg=DarkText}
\setbeamercolor{normal text}{fg=DarkText}
\setbeamercolor{alerted text}{fg=RedTitle}
\setbeamersize{text margin left=5mm, text margin right=5mm}
\addtobeamertemplate{frametitle}{}{\vspace{-0.5em}}

\setbeamertemplate{navigation symbols}{}
\setbeamertemplate{footline}{%
  \hspace{5mm}{\usebeamerfont{page number in head/foot}\color{black!45}%
  ECON 282E \,$\cdot$\, \insertshortdeck}\hfill%
  \usebeamercolor[fg]{page number in head/foot}%
  \usebeamerfont{page number in head/foot}%
  \insertframenumber\,/\,\inserttotalframenumber\kern1em\vskip2pt%
}

% ---------------------------------------------------------------------------
% Chinese text under pdflatex (book titles etc.)
\newcommand{\zh}[1]{\begin{CJK*}{UTF8}{gbsn}#1\end{CJK*}}

% ---------------------------------------------------------------------------
% Deck title block
%   \decktitle{session}{deck id}{title}{date}
\newcommand{\insertshortdeck}{}
\newcommand{\decksession}{}
\newcommand{\deckid}{}
\newcommand{\decktitletext}{}
\newcommand{\deckdate}{}
\newcommand{\decktitle}[4]{%
  \renewcommand{\decksession}{#1}%
  \renewcommand{\deckid}{#2}%
  \renewcommand{\decktitletext}{#3}%
  \renewcommand{\deckdate}{#4}%
  \renewcommand{\insertshortdeck}{Session #1 \,$\cdot$\, Deck #1.#2}%
  \title{#3}%
  \author{Zhigang Feng}%
  \date{#4}%
}
\newcommand{\makedecktitle}{%
  \begin{frame}[plain,noframenumbering]
    \begin{tikzpicture}[remember picture,overlay]
      \fill[UCRBlue] (current page.north west) rectangle ([yshift=-0.55cm]current page.north east);
      \fill[UCRGold] ([yshift=-0.55cm]current page.north west) rectangle ([yshift=-0.62cm]current page.north east);
      \node[anchor=west, text=white, font=\small\bfseries] at ([xshift=5mm,yshift=-0.275cm]current page.north west)
        {ECON 282E \,$\cdot$\, Foundations of Macroeconomics};
      \node[anchor=east, text=white, font=\small] at ([xshift=-5mm,yshift=-0.275cm]current page.north east)
        {UC Riverside \,$\cdot$\, Fall 2026};
    \end{tikzpicture}
    \vspace*{1.1cm}
    {\usebeamerfont{title}\color{black!55}\large Session \decksession\ \,$\cdot$\, Deck \decksession.\deckid\par}
    \vspace{0.25cm}
    {\usebeamerfont{title}\color{RedTitle}\LARGE\bfseries \decktitletext\par}
    \vspace{0.7cm}
    {\large Zhigang Feng\par}
    \vspace{0.15cm}
    {\color{black!65}\deckdate\par}
    \vfill
    {\footnotesize\itshape\color{black!55}Build it on paper $\;\to\;$ solve it on the machine $\;\to\;$ push it to the frontier with AI\par}
  \end{frame}}

% ---------------------------------------------------------------------------
% Section divider (plain frame, not counted)
\newcommand{\sectiondivider}[2]{%
\begin{frame}[plain,noframenumbering]
\begin{center}
\vfill
{\Large\textbf{#1}\par}
\vspace{0.35cm}
{\normalsize #2\par}
\vfill
\end{center}
\end{frame}
}

% ---------------------------------------------------------------------------
% Boxes
\newtcolorbox{conceptbox}[1][]{colback=blue!3, colframe=RedTitle!55, boxrule=0.35mm,
  arc=1mm, left=2mm, right=2mm, top=1mm, bottom=1mm, #1}
\newtcolorbox{cautionbox}[1][]{colback=red!3, colframe=red!50!black, boxrule=0.35mm,
  arc=1mm, left=2mm, right=2mm, top=1mm, bottom=1mm, #1}
\newtcolorbox{examplebox}[1][]{colback=green!3, colframe=green!45!black, boxrule=0.35mm,
  arc=1mm, left=2mm, right=2mm, top=1mm, bottom=1mm, #1}
\newtcolorbox{takeawaybox}[1][]{colback=VeryLightGray, colframe=RedTitle!55, boxrule=0.35mm,
  arc=1mm, left=2mm, right=2mm, top=1mm, bottom=1mm, #1}
\newtcolorbox{researchbox}[1][]{enhanced, colback=violet!4, colframe=violet!60!black,
  boxrule=0.35mm, arc=1mm, left=2mm, right=2mm, top=1mm, bottom=1mm,
  fonttitle=\bfseries\small, title={Research in the wild}, #1}

\newcommand{\compacttables}{\renewcommand{\arraystretch}{1.15}}
\newcommand{\src}[1]{{\tiny\color{black!55}#1}}

% ---------------------------------------------------------------------------
%% Course map: the recurring "you are here" slide for ECON 282E.
%% Loaded by course_preamble.tex. Pattern follows AI-ECON-2026 lec01_map.tex, but the map
%% shows the whole ten-session course rather than one lecture.
%%
%%   \CourseMap{s}{label}   s = 1..10 highlights session s and prints "you are here: label"
%%                          (e.g. \CourseMap{1}{1.A2}); earlier sessions get a check mark.
%%   \CourseMap{0}{}        full overview, nothing highlighted.
%%
%% Note: TikZ option lists are comma-split before TeX conditionals run, so every \ifnum
%% below sits outside [...] option lists.

\newcommand{\CourseMap}[2]{%
\begin{center}
\begin{tikzpicture}[
    sess/.style={rectangle, rounded corners=2pt, draw=black!45, minimum width=1.34cm,
        minimum height=1.12cm, text width=1.26cm, align=center, inner sep=1pt},
    stage/.style={font=\tiny\bfseries, text=black!70},
]
  \foreach \i/\d/\t/\c in {%
      1/Sep 24/Intro \\ Growth/StageTrunk,
      2/Oct 1/HA $\cdot$ OLG \\ Policy/StageBranch,
      3/Oct 8/Num.\ DP \\ Python/StageMachine,
      4/Oct 15/Numerics \\ I/StageMachine,
      5/Oct 22/Numerics \\ II/StageMachine,
      6/Oct 29/The wall \\ \textbf{Midterm}/StageWall,
      7/Nov 5/ML \\ Deep nets/StageTools,
      8/Nov 12/RL \\ HA $+$ DL/StageTools,
      9/Nov 19/Text \\ LLMs/StageData,
      10/Dec 3/Agents \\ \textbf{Projects}/StageAgents} {
    \node[sess, fill=\c] (n\i) at ({1.46*(\i-1)}, 0)
      {{\scriptsize\bfseries S\i}\\[-1pt]{\tiny\color{black!60}\d}\\[1pt]{\tiny \t}};
    \ifnum\i<#1
      \node[font=\scriptsize, text=green!45!black] at ([xshift=-2.5mm,yshift=-2.2mm]n\i.north east) {$\checkmark$};
    \fi
    \ifnum\i=#1
      \draw[RedTitle, very thick, rounded corners=2pt]
        ([xshift=-0.6mm,yshift=-0.6mm]n\i.south west) rectangle ([xshift=0.6mm,yshift=0.6mm]n\i.north east);
      % keep the label inside the picture at both ends of the row
      \ifnum\i<3
        \node[font=\scriptsize\bfseries, text=RedTitle, anchor=south west] at ([yshift=1.2mm]n\i.north west)
          {$\blacktriangledown$\,you are here: #2};
      \else\ifnum\i>8
        \node[font=\scriptsize\bfseries, text=RedTitle, anchor=south east] at ([yshift=1.2mm]n\i.north east)
          {you are here: #2\,$\blacktriangledown$};
      \else
        \node[font=\scriptsize\bfseries, text=RedTitle, anchor=south] at ([yshift=1.2mm]n\i.north)
          {you are here: #2\,$\blacktriangledown$};
      \fi\fi
    \fi
  }
  % stage band under the sessions
  \foreach \a/\b/\lab/\c in {1/1/trunk/StageTrunk, 2/2/branches/StageBranch,
      3/5/the machine $\cdot$ classical methods/StageMachine, 6/6/the wall/StageWall,
      7/8/new tools/StageTools, 9/9/new data/StageData, 10/10/colleagues/StageAgents} {
    \fill[\c] ([yshift=-1.5mm]n\a.south west) rectangle ([yshift=-4.6mm]n\b.south east);
    \path ([yshift=-1.5mm]n\a.south west) -- ([yshift=-4.6mm]n\b.south east)
      node[midway, stage] {\lab};
  }
  \node[font=\footnotesize\itshape, text=black!70] at ([yshift=-9.5mm]$(n1.south)!0.5!(n10.south)$)
    {Build it on paper $\;\to\;$ solve it on the machine $\;\to\;$ push it to the frontier with AI};
\end{tikzpicture}
\end{center}
}


%%   \decktitle{1}{A1}{Who I Am \& Why This Course}{September 24, 2026}
%%   \begin{document}
%%   \makedecktitle
%%   ...
%%
%% Adapted from Courses/AI-ECON-2026/source/shared_preamble.tex (Metropolis theme,
%% concept/caution/example/takeaway boxes, \sectiondivider). Changes for this course:
%%   * course title block (\decktitle / \makedecktitle) instead of \makebeamertitle
%%   * \zh{...} typesets Chinese (book titles) under pdflatex via CJKutf8
%%   * researchbox for the "Research in the wild" frames
%%   * section pages off; TikZ libraries and the course map loaded here
%% Build with pdflatex only (two passes); tools/build_slides.py does this for you.

\documentclass[10pt,english,aspectratio=169]{beamer}

\usepackage[T1]{fontenc}
\usepackage{textcomp}
\usepackage[utf8]{inputenc}
\usepackage{babel}
\usepackage{CJKutf8}

\usepackage{amsmath, amssymb, amsthm, graphicx}
\usepackage{booktabs, multirow, tabularx}
\usepackage{tikz}
\usetikzlibrary{positioning,arrows.meta,shapes,calc,fit,backgrounds}
\usepackage{pgfplots}
\pgfplotsset{compat=1.16}
\usepackage[most]{tcolorbox}
\tcbuselibrary{skins, breakable}
\usepackage{listings}
\usepackage{xcolor}

\lstset{
  basicstyle=\ttfamily\small,
  keywordstyle=\color{blue!80!black}\bfseries,
  commentstyle=\color{green!50!black}\itshape,
  stringstyle=\color{orange!80!black},
  backgroundcolor=\color{gray!8},
  frame=single,
  rulecolor=\color{gray!40},
  breaklines=true,
  showstringspaces=false,
  numbers=none,
}

\usetheme{metropolis}
\usefonttheme{professionalfonts}
\metroset{sectionpage=none}

\definecolor{LightGray}{RGB}{230,230,230}
\definecolor{RedTitle}{RGB}{0,0,200}
\definecolor{VeryLightGray}{RGB}{245,245,245}
\definecolor{DarkText}{RGB}{0,0,0}
\definecolor{UCRBlue}{RGB}{0,45,106}
\definecolor{UCRGold}{RGB}{241,171,0}

% Stage colours of the course arc (used by the course map and elsewhere)
\colorlet{StageTrunk}{blue!12}
\colorlet{StageBranch}{cyan!14}
\colorlet{StageMachine}{gray!18}
\colorlet{StageWall}{red!14}
\colorlet{StageTools}{orange!20}
\colorlet{StageData}{green!16}
\colorlet{StageAgents}{violet!14}

\hypersetup{bookmarks=false,breaklinks=true,pdfborder={0 0 0},colorlinks=true,
  linkcolor=DarkText,urlcolor=RedTitle}

\setbeamercolor{background canvas}{bg=white}
\setbeamercolor{title}{fg=RedTitle}
\setbeamercolor{frametitle}{bg=VeryLightGray,fg=DarkText}
\setbeamercolor{normal text}{fg=DarkText}
\setbeamercolor{alerted text}{fg=RedTitle}
\setbeamersize{text margin left=5mm, text margin right=5mm}
\addtobeamertemplate{frametitle}{}{\vspace{-0.5em}}

\setbeamertemplate{navigation symbols}{}
\setbeamertemplate{footline}{%
  \hspace{5mm}{\usebeamerfont{page number in head/foot}\color{black!45}%
  ECON 282E \,$\cdot$\, \insertshortdeck}\hfill%
  \usebeamercolor[fg]{page number in head/foot}%
  \usebeamerfont{page number in head/foot}%
  \insertframenumber\,/\,\inserttotalframenumber\kern1em\vskip2pt%
}

% ---------------------------------------------------------------------------
% Chinese text under pdflatex (book titles etc.)
\newcommand{\zh}[1]{\begin{CJK*}{UTF8}{gbsn}#1\end{CJK*}}

% ---------------------------------------------------------------------------
% Deck title block
%   \decktitle{session}{deck id}{title}{date}
\newcommand{\insertshortdeck}{}
\newcommand{\decksession}{}
\newcommand{\deckid}{}
\newcommand{\decktitletext}{}
\newcommand{\deckdate}{}
\newcommand{\decktitle}[4]{%
  \renewcommand{\decksession}{#1}%
  \renewcommand{\deckid}{#2}%
  \renewcommand{\decktitletext}{#3}%
  \renewcommand{\deckdate}{#4}%
  \renewcommand{\insertshortdeck}{Session #1 \,$\cdot$\, Deck #1.#2}%
  \title{#3}%
  \author{Zhigang Feng}%
  \date{#4}%
}
\newcommand{\makedecktitle}{%
  \begin{frame}[plain,noframenumbering]
    \begin{tikzpicture}[remember picture,overlay]
      \fill[UCRBlue] (current page.north west) rectangle ([yshift=-0.55cm]current page.north east);
      \fill[UCRGold] ([yshift=-0.55cm]current page.north west) rectangle ([yshift=-0.62cm]current page.north east);
      \node[anchor=west, text=white, font=\small\bfseries] at ([xshift=5mm,yshift=-0.275cm]current page.north west)
        {ECON 282E \,$\cdot$\, Foundations of Macroeconomics};
      \node[anchor=east, text=white, font=\small] at ([xshift=-5mm,yshift=-0.275cm]current page.north east)
        {UC Riverside \,$\cdot$\, Fall 2026};
    \end{tikzpicture}
    \vspace*{1.1cm}
    {\usebeamerfont{title}\color{black!55}\large Session \decksession\ \,$\cdot$\, Deck \decksession.\deckid\par}
    \vspace{0.25cm}
    {\usebeamerfont{title}\color{RedTitle}\LARGE\bfseries \decktitletext\par}
    \vspace{0.7cm}
    {\large Zhigang Feng\par}
    \vspace{0.15cm}
    {\color{black!65}\deckdate\par}
    \vfill
    {\footnotesize\itshape\color{black!55}Build it on paper $\;\to\;$ solve it on the machine $\;\to\;$ push it to the frontier with AI\par}
  \end{frame}}

% ---------------------------------------------------------------------------
% Section divider (plain frame, not counted)
\newcommand{\sectiondivider}[2]{%
\begin{frame}[plain,noframenumbering]
\begin{center}
\vfill
{\Large\textbf{#1}\par}
\vspace{0.35cm}
{\normalsize #2\par}
\vfill
\end{center}
\end{frame}
}

% ---------------------------------------------------------------------------
% Boxes
\newtcolorbox{conceptbox}[1][]{colback=blue!3, colframe=RedTitle!55, boxrule=0.35mm,
  arc=1mm, left=2mm, right=2mm, top=1mm, bottom=1mm, #1}
\newtcolorbox{cautionbox}[1][]{colback=red!3, colframe=red!50!black, boxrule=0.35mm,
  arc=1mm, left=2mm, right=2mm, top=1mm, bottom=1mm, #1}
\newtcolorbox{examplebox}[1][]{colback=green!3, colframe=green!45!black, boxrule=0.35mm,
  arc=1mm, left=2mm, right=2mm, top=1mm, bottom=1mm, #1}
\newtcolorbox{takeawaybox}[1][]{colback=VeryLightGray, colframe=RedTitle!55, boxrule=0.35mm,
  arc=1mm, left=2mm, right=2mm, top=1mm, bottom=1mm, #1}
\newtcolorbox{researchbox}[1][]{enhanced, colback=violet!4, colframe=violet!60!black,
  boxrule=0.35mm, arc=1mm, left=2mm, right=2mm, top=1mm, bottom=1mm,
  fonttitle=\bfseries\small, title={Research in the wild}, #1}

\newcommand{\compacttables}{\renewcommand{\arraystretch}{1.15}}
\newcommand{\src}[1]{{\tiny\color{black!55}#1}}

% ---------------------------------------------------------------------------
%% Course map: the recurring "you are here" slide for ECON 282E.
%% Loaded by course_preamble.tex. Pattern follows AI-ECON-2026 lec01_map.tex, but the map
%% shows the whole ten-session course rather than one lecture.
%%
%%   \CourseMap{s}{label}   s = 1..10 highlights session s and prints "you are here: label"
%%                          (e.g. \CourseMap{1}{1.A2}); earlier sessions get a check mark.
%%   \CourseMap{0}{}        full overview, nothing highlighted.
%%
%% Note: TikZ option lists are comma-split before TeX conditionals run, so every \ifnum
%% below sits outside [...] option lists.

\newcommand{\CourseMap}[2]{%
\begin{center}
\begin{tikzpicture}[
    sess/.style={rectangle, rounded corners=2pt, draw=black!45, minimum width=1.34cm,
        minimum height=1.12cm, text width=1.26cm, align=center, inner sep=1pt},
    stage/.style={font=\tiny\bfseries, text=black!70},
]
  \foreach \i/\d/\t/\c in {%
      1/Sep 24/Intro \\ Growth/StageTrunk,
      2/Oct 1/HA $\cdot$ OLG \\ Policy/StageBranch,
      3/Oct 8/Num.\ DP \\ Python/StageMachine,
      4/Oct 15/Numerics \\ I/StageMachine,
      5/Oct 22/Numerics \\ II/StageMachine,
      6/Oct 29/The wall \\ \textbf{Midterm}/StageWall,
      7/Nov 5/ML \\ Deep nets/StageTools,
      8/Nov 12/RL \\ HA $+$ DL/StageTools,
      9/Nov 19/Text \\ LLMs/StageData,
      10/Dec 3/Agents \\ \textbf{Projects}/StageAgents} {
    \node[sess, fill=\c] (n\i) at ({1.46*(\i-1)}, 0)
      {{\scriptsize\bfseries S\i}\\[-1pt]{\tiny\color{black!60}\d}\\[1pt]{\tiny \t}};
    \ifnum\i<#1
      \node[font=\scriptsize, text=green!45!black] at ([xshift=-2.5mm,yshift=-2.2mm]n\i.north east) {$\checkmark$};
    \fi
    \ifnum\i=#1
      \draw[RedTitle, very thick, rounded corners=2pt]
        ([xshift=-0.6mm,yshift=-0.6mm]n\i.south west) rectangle ([xshift=0.6mm,yshift=0.6mm]n\i.north east);
      % keep the label inside the picture at both ends of the row
      \ifnum\i<3
        \node[font=\scriptsize\bfseries, text=RedTitle, anchor=south west] at ([yshift=1.2mm]n\i.north west)
          {$\blacktriangledown$\,you are here: #2};
      \else\ifnum\i>8
        \node[font=\scriptsize\bfseries, text=RedTitle, anchor=south east] at ([yshift=1.2mm]n\i.north east)
          {you are here: #2\,$\blacktriangledown$};
      \else
        \node[font=\scriptsize\bfseries, text=RedTitle, anchor=south] at ([yshift=1.2mm]n\i.north)
          {you are here: #2\,$\blacktriangledown$};
      \fi\fi
    \fi
  }
  % stage band under the sessions
  \foreach \a/\b/\lab/\c in {1/1/trunk/StageTrunk, 2/2/branches/StageBranch,
      3/5/the machine $\cdot$ classical methods/StageMachine, 6/6/the wall/StageWall,
      7/8/new tools/StageTools, 9/9/new data/StageData, 10/10/colleagues/StageAgents} {
    \fill[\c] ([yshift=-1.5mm]n\a.south west) rectangle ([yshift=-4.6mm]n\b.south east);
    \path ([yshift=-1.5mm]n\a.south west) -- ([yshift=-4.6mm]n\b.south east)
      node[midway, stage] {\lab};
  }
  \node[font=\footnotesize\itshape, text=black!70] at ([yshift=-9.5mm]$(n1.south)!0.5!(n10.south)$)
    {Build it on paper $\;\to\;$ solve it on the machine $\;\to\;$ push it to the frontier with AI};
\end{tikzpicture}
\end{center}
}


\graphicspath{{./figures/}}

%% Block B of this session is the first user of the preamble's listings setup,
%% which sets no language.  Fix that locally rather than touching the shared file.
\lstset{language=Python, basicstyle=\ttfamily\scriptsize, frame=none,
        backgroundcolor=\color{gray!8}, aboveskip=2pt, belowskip=2pt}

\decktitle{3}{A1}{What Does It Mean to Solve a Model?}{Thursday, October 8, 2026}

\begin{document}

\makedecktitle

%=============================================================================
\begin{frame}{Where We Are}
\CourseMap{3}{3.A1}
\vspace{-1mm}
\begin{center}
\small \textbf{Block A:} \textbf{3.A1} what it means to solve a model $\;\cdot\;$ \textbf{3.A2} value function iteration on a grid $\;\cdot\;$ \textbf{3.A3} Euler equations and time iteration $\;\cdot\;$ \textbf{3.A4} the worked example, end to end.
\end{center}
\end{frame}

%=============================================================================
\begin{frame}[fragile]{We Proved It Exists. Now Produce It.}
\begin{columns}[T]
\begin{column}{0.53\textwidth}
\small
Two weeks of theorems, and every one of them is an \textbf{existence} statement:
\begin{itemize}\footnotesize
  \item $T$ is a contraction, so a unique $v^*$ exists and $T^n v_0 \to v^*$.
  \item A unique invariant $\lambda^*$ exists, and $\lambda_0 Q^n \Rightarrow \lambda^*$ from any start.
  \item Capital supply and demand cross, so an $r^*$ exists.
\end{itemize}
\medskip
Not one of them is a \textbf{number}. None tells you what $v^*(0.3)$ is, how unequal $\lambda^*$ is, or whether $r^*$ is 2\% or 5\%.
\medskip

\textbf{Today: what stands between the theorem and the number} --- and what you must learn in order to cross it.
\end{column}
\begin{column}{0.44\textwidth}
\begin{lstlisting}
>>> 0.1 + 0.2
0.30000000000000004
>>> 0.1 + 0.2 == 0.3
False
\end{lstlisting}
\vspace{-1mm}
\begin{cautionbox}\footnotesize
We are about to ask a machine for a fixed point in an infinite-dimensional function space. It cannot represent one tenth.
\end{cautionbox}
\end{column}
\end{columns}
\end{frame}

%=============================================================================
\begin{frame}{The Unknown Is a Function}
\begin{columns}[T]
\begin{column}{0.54\textwidth}
\small
In a static problem you solve for numbers. Here the Bellman equation
\[
  v(k) = \max_{k'\in\Gamma(k)}\big\{u\big(zf(k)-k'\big) + \beta\,\mathbb{E}\,v(k')\big\}
\]
is one equation for an \textbf{entire function} $v$. Bellman's move was to replace an infinite sequence of choices with a \textbf{functional equation} --- a fixed point in a function space.
\medskip

Same on the Euler road. Substituting $k'=g(k)$ into $u'(c)=\beta u'(c')R$ gives
\[
  u'\big(zf(k)-g(k)\big)=\beta\,u'\big(zf(g(k))-g(g(k))\big)R\big(g(k)\big),
\]
where the unknown $g$ appears \emph{composed with itself}.
\end{column}
\begin{column}{0.43\textwidth}
\begin{conceptbox}[title=\textbf{The gap, stated once}]\small
$v$ and $g$ live in an infinite-dimensional space. Your computer holds finitely many numbers, each of them approximate.
\end{conceptbox}
\medskip
\begin{examplebox}\footnotesize
So every method in this course does the same two things: pick a \textbf{finite description} of a function, and pick a \textbf{finite procedure} that drives the residual to zero.
\end{examplebox}
\vspace{1mm}
\src{Bellman (1957); the composed-policy form follows Judd (1998), ch.~12.}
\end{column}
\end{columns}
\end{frame}

%=============================================================================
\begin{frame}{Two Roads --- and Both Are Constructive}
\renewcommand{\arraystretch}{1.2}
\begin{center}\footnotesize
\begin{tabularx}{\linewidth}{>{\raggedright\arraybackslash}p{2.5cm} >{\raggedright\arraybackslash}X >{\raggedright\arraybackslash}X}
\toprule
 & \textbf{Road 1: dynamic programming} & \textbf{Road 2: Euler equations} \\
\midrule
Unknown & the value function $v$ & the policy function $g$ \\
Iteration & $v_{n+1}=Tv_n$ & $g_{n+1}=Kg_n$ (one root-find per point) \\
Why it converges & $T$ is a \textbf{contraction} (Blackwell) & $K$ is \textbf{monotone} and bounded \\
Needs & only the Bellman equation & differentiability, interiority \\
\bottomrule
\end{tabularx}
\end{center}
\vspace{1mm}
\begin{takeawaybox}\small
The theorems of 1.B2 are not decoration. A contraction argument is an \textbf{algorithm with an error bound attached}; a monotonicity argument is an algorithm with a convergence guarantee. Both roads are constructive --- that is precisely why we could prove existence in the first place.
\end{takeawaybox}
\vspace{1mm}
\src{Adapted from QM ``Introduction to Quantitative Macroeconomics'', ``Two Roads to Equilibrium''. Road 2 is Coleman (1990) and Greenwood--Huffman (1995); we build it in 3.A3.}
\end{frame}

%=============================================================================
\begin{frame}{Six Decisions Before You Can Type Anything}
\renewcommand{\arraystretch}{1.15}
\begin{center}\footnotesize
\begin{tabularx}{\linewidth}{@{}>{\raggedright\arraybackslash}p{0.5cm} >{\raggedright\arraybackslash}p{3.3cm} >{\raggedright\arraybackslash}X >{\raggedright\arraybackslash}p{3.0cm}@{}}
\toprule
 & \textbf{Decision} & \textbf{Question it answers} & \textbf{Today's crude answer} \\
\midrule
1 & Represent the function & How is $v$ stored in finite memory? & values on a grid \\
2 & Discretize the state & Which points, how many, spaced how? & uniform grid \\
3 & Evaluate off the grid & What is $v(k')$ between two points? & linear interpolation \\
4 & Handle the choice & How is the $\max$ (or the root) taken? & search the grid \\
5 & Take the expectation & How does $\mathbb{E}$ become a finite sum? & a \emph{given} Markov chain \\
6 & Iterate and stop & When is it converged, and how wrong is it? & $\lVert v_{n+1}-v_n\rVert<\varepsilon$ \\
\bottomrule
\end{tabularx}
\end{center}
\vspace{1mm}
\begin{conceptbox}\small
This is the list. \textbf{Every numerical method in this course is one choice from each row} --- and the whole of Sessions 4 and 5 is about replacing the right-hand column with something better.
\end{conceptbox}
\end{frame}

%=============================================================================
\begin{frame}{Three Sources of Error --- and a Fourth}
\begin{columns}[T]
\begin{column}{0.52\textwidth}
\small
Replacing the functional equation by something finite introduces error in three places at once:
\begin{enumerate}\footnotesize
  \item \textbf{State discretization.} $v$ is known at $N$ points; everywhere else it is a guess.
  \item \textbf{Expectation approximation.} A continuous shock becomes a finite chain; $\int$ becomes $\sum$.
  \item \textbf{Choice-set restriction.} The maximizer is searched over a finite set, not over $\Gamma(k)$.
\end{enumerate}
\medskip
Each is a modelling decision you made, and each can be measured.
\end{column}
\begin{column}{0.45\textwidth}
\begin{cautionbox}[title=\textbf{The fourth one is not in the theory}]\footnotesize
\textbf{Arithmetic.} Every one of those finite operations is carried out in floating point. The numbers your machine adds are not the numbers you wrote down.
\end{cautionbox}
\medskip
\begin{examplebox}\footnotesize
It is the only error source you cannot reduce by thinking harder about economics --- so it is worth ten minutes now, once, for the rest of your career.
\end{examplebox}
\end{column}
\end{columns}
\vspace{1mm}
\src{Sources 1--3 adapted from QM Lec.~7, ``VFI: From Continuous to Discrete''.}
\end{frame}

%=============================================================================
\begin{frame}{Floating Point: How a Real Number Is Actually Stored}
\begin{columns}[T]
\begin{column}{0.54\textwidth}
\small
A double is 64 bits, read as
\[
  (-1)^{S}\,\underbrace{(1+M)}_{\text{mantissa}}\,2^{\,E-1023},
\]
with $S$ one sign bit, $E$ an 11-bit exponent \emph{field} and $M$ a 52-bit fraction; 1023 is a \textbf{bias}, so one unsigned field encodes both large and small exponents.
\medskip

\textbf{Read one off.} $6.5$ is stored as
\begin{center}\ttfamily\scriptsize
0\;\;10000000001\;\;1010\ldots0
\end{center}
\footnotesize
$S=0$ (positive) $\cdot$ $E=1025$, so the exponent is $1025-1023=2$ $\cdot$ $M$ gives a mantissa of $1.625$ $\cdot$ and $1.625\times2^{2}=\mathbf{6.5}$.
\end{column}
\begin{column}{0.43\textwidth}
\begin{center}\footnotesize
\begin{tabular}{@{}lr@{}}
\toprule
\multicolumn{2}{@{}l}{\textbf{The four boundaries (float64)}} \\
\midrule
largest & $1.80\times10^{308}$ \\
machine epsilon & $2.22\times10^{-16}$ \\
smallest normal & $2.23\times10^{-308}$ \\
smallest subnormal & $4.94\times10^{-324}$ \\
\bottomrule
\end{tabular}
\end{center}
\medskip
\begin{conceptbox}\footnotesize
Between the last two lies \textbf{subnormal} territory: below $2^{-1022}$ the leading 1 is surrendered to buy a few more powers of two, and precision degrades as you go.
\end{conceptbox}
\vspace{1mm}
{\footnotesize The exponent buys \textbf{range}; only the mantissa buys \textbf{precision}.}
\end{column}
\end{columns}
\vspace{1mm}
\src{Adapted from QM ``Numerical Methods'', floating-point block; the source works the example in 32-bit, this course runs float64. Bit pattern read with \texttt{struct}.}
\end{frame}

%=============================================================================
\begin{frame}{Floating Point: The Real Line Is Not on Your Computer}
\begin{columns}[T]
\begin{column}{0.44\textwidth}
\small
A double is stored as
\[
  (-1)^{S}\,(1+M)\,2^{E-1023},
\]
with $1$ sign bit, $11$ exponent bits and $52$ mantissa bits. The mantissa buys precision; the exponent only slides you along the line.
\medskip

So the representable numbers are \textbf{not evenly spaced}. Between $2^{j}$ and $2^{j+1}$ they sit $2^{j}\varepsilon_{\text{mach}}$ apart --- the gap \textbf{doubles at every power of two}.
\medskip
\begin{center}\footnotesize
\begin{tabular}{@{}lll@{}}
\toprule
 & $\varepsilon_{\text{mach}}$ & digits \\
\midrule
float64 & $2^{-52}\approx2.2\times10^{-16}$ & $\approx16$ \\
float32 & $2^{-23}\approx1.2\times10^{-7}$ & $\approx7$ \\
\bottomrule
\end{tabular}
\end{center}
\end{column}
\begin{column}{0.53\textwidth}
\centering
\begin{tikzpicture}[font=\scriptsize, yscale=1.0]
\draw[->, thick, black!60] (-0.15,0) -- (7.35,0);
\foreach \i in {0,...,8} \draw[RedTitle, thick]        (\i*0.225,-0.13) -- (\i*0.225,0.13);
\foreach \i in {0,...,8} \draw[orange!80!black, thick] (1.8+\i*0.45,-0.13) -- (1.8+\i*0.45,0.13);
\foreach \i in {0,...,2} \draw[black!50, thick]        (5.4+\i*0.9,-0.13) -- (5.4+\i*0.9,0.13);
\foreach \x/\l in {0/$1$, 1.8/$2$, 5.4/$4$}
  \node[font=\scriptsize, anchor=north, text=black!70] at (\x,-0.2) {\l};
\node[font=\tiny, text=RedTitle, anchor=south]        at (0.9,0.22) {gap $\varepsilon$};
\node[font=\tiny, text=orange!80!black, anchor=south] at (3.6,0.22) {gap $2\varepsilon$};
\node[font=\tiny, text=black!50, anchor=south]        at (6.3,0.22) {gap $4\varepsilon$};
\node[font=\scriptsize, anchor=west, text=black!60]   at (7.0,-0.28) {$\dots$};
\end{tikzpicture}
\\[3mm]
\begin{cautionbox}\footnotesize
$\varepsilon_{\text{mach}}$ is a \textbf{relative} precision. There is no fixed number of decimal places --- only a fixed number of significant digits, wherever on the line you happen to be standing.
\end{cautionbox}
\end{column}
\end{columns}
\vspace{0.5mm}
\src{Adapted from QM ``Numerical Methods'', floating-point section; the 32-bit version there makes the same point with $2^{-23}$.}
\end{frame}

%=============================================================================
\begin{frame}{What That Does to a Macro Model}
\begin{columns}[T]
\begin{column}{0.50\textwidth}
\small
Macro models routinely hold quantities that differ by fifteen orders of magnitude in the same array: $K$ and $Y$ around $10^{12}$--$10^{14}$, while $r$, $\delta$ and a convergence tolerance live at $10^{-2}$--$10^{-8}$.
\medskip

At US GDP $\approx\$28$ trillion, the distance to the next representable number is
\begin{center}\footnotesize
\begin{tabular}{@{}lr@{}}
\toprule
\textbf{float64} & \$0.0039 \\
\textbf{float32} & \$2{,}097{,}152 \\
\bottomrule
\end{tabular}
\end{center}
\medskip
In single precision, \textbf{adding \$100{,}000 to \$28 trillion changes nothing at all}. The addition is silently discarded, and no error is raised.
\end{column}
\begin{column}{0.47\textwidth}
\begin{cautionbox}[title=\textbf{The ``big $+$ small'' trap}]\footnotesize
Whenever a small correction is added to a large accumulator --- a market-clearing residual, a welfare sum, a capital stock --- check that the small thing is not being absorbed.
\end{cautionbox}
\medskip
\begin{takeawaybox}\footnotesize
\textbf{Course default: float64.} Machine learning defaults to float32 (or less) because gradients tolerate noise. \emph{Accounting identities do not.} When we reach PyTorch in 3.B3, this is the first setting to check.
\end{takeawaybox}
\end{column}
\end{columns}
\vspace{1mm}
\src{Spacings computed with \texttt{numpy.spacing}; all values recorded in the session's number ledger.}
\end{frame}

%=============================================================================
\begin{frame}{Three Ways the Arithmetic Lies}
\begin{columns}[T]
\begin{column}{0.51\textwidth}
\small
\textbf{1. Cancellation.} Subtracting nearly equal numbers destroys the leading digits, and what is left is noise promoted to the front:
\begin{center}\ttfamily\scriptsize
(1.0 + 1e-16) - 1.0 \;$\to$\; 0.0
\end{center}
\textbf{2. Addition is not associative.}
\begin{center}\ttfamily\scriptsize
(0.1+0.2)+0.3 $\to$ 0.6000000000000001\\
0.1+(0.2+0.3) $\to$ 0.6
\end{center}
so the \emph{order} of a sum changes its value --- and a parallel reduction may not even fix the order between runs.
\medskip

\textbf{3. Overflow, underflow, and \texttt{nan}.} $10^{308}\!\times\!10\to$ \texttt{inf}; $10^{-320}/10^{10}\to$ \texttt{0.0}; and \texttt{nan != nan} is \texttt{True}.
\end{column}
\begin{column}{0.46\textwidth}
\begin{cautionbox}[title=\textbf{Where each one finds you}]\footnotesize
\textbf{Cancellation}: a finite-difference derivative gets \emph{worse} as the step shrinks, because the numerator is a difference of nearly equal numbers (measured in 3.B3).\\[1mm]
\textbf{Ordering}: a market-clearing residual moves in the last digits when you change how you sum.\\[1mm]
\textbf{\texttt{nan}} loses every comparison silently --- one of them can make \texttt{max} return a wrong answer, unflagged.
\end{cautionbox}
\medskip
\begin{takeawaybox}\footnotesize
Practical consequences: mark infeasible choices with a \textbf{large negative number, not $-\infty$}; compute a residual as a \textbf{ratio}, not a difference; and never test floats with \texttt{==}.
\end{takeawaybox}
\end{column}
\end{columns}
\end{frame}

%=============================================================================
\begin{frame}{Three Error Budgets, and How They Compose}
\begin{columns}[T]
\begin{column}{0.54\textwidth}
\small
Your answer differs from the truth for three independent reasons:
\medskip
\begin{center}\footnotesize
\begin{tabularx}{\linewidth}{@{}>{\raggedright\arraybackslash}p{2.1cm} >{\raggedright\arraybackslash}X >{\raggedright\arraybackslash}p{1.9cm}@{}}
\toprule
\textbf{Budget} & \textbf{Comes from} & \textbf{Controlled by} \\
\midrule
Approximation & finitely many grid points & $h$, the spacing \\
Truncation & stopping after $n$ iterations & $\varepsilon$, the tolerance \\
Round-off & floating-point arithmetic & the dtype \\
\bottomrule
\end{tabularx}
\end{center}
\medskip
They \textbf{add}. The total is no smaller than the largest of the three, so paying to shrink one below the others buys you exactly nothing.
\end{column}
\begin{column}{0.43\textwidth}
\begin{conceptbox}[title=\textbf{A rule you can apply today}]\small
Before refining a grid, ask what your iteration tolerance already costs you. Before tightening a tolerance, ask what your grid already costs you.
\end{conceptbox}
\medskip
\begin{examplebox}\footnotesize
In 3.A2 we will watch this happen: a grid refined from 400 to 800 points delivers \textbf{no improvement at all} in the value function, because the iteration tolerance has already put a floor underneath it.
\end{examplebox}
\end{column}
\end{columns}
\end{frame}

%=============================================================================
\begin{frame}{What ``Right'' Means: A Number Is a Claim}
\begin{columns}[T]
\begin{column}{0.50\textwidth}
\small
You almost never know the truth, so accuracy is reported as a \textbf{residual}: how badly does the computed policy violate an equilibrium condition it was never given?
\[
  \mathcal{E}(k,z)=\bigg|1-\frac{\tilde c(k,z)}{c(k,z)}\bigg|,
\]
the \textbf{unit-free} consumption mistake of 1.B2. It reads directly as money: $\log_{10}\mathcal{E}=-3$ is a one-dollar error per \$1{,}000 of consumption.
\vspace{2mm}
\begin{cautionbox}[title=\textbf{``A decade'' --- the unit we report in}]\footnotesize
Because errors span ten orders of magnitude, we compare them on the $\log_{10}$ scale and call one unit a \textbf{decade}: a \textbf{factor of ten} in the error. $10^{-2}$ to $10^{-3}$ is one decade --- one digit bought. ``1.2 decades'' means the error fell by $10^{1.2}\approx16$ times.
\end{cautionbox}
\end{column}
\begin{column}{0.47\textwidth}
\begin{center}\footnotesize
\begin{tabular}{@{}lll@{}}
\toprule
$\log_{10}\mathcal{E}$ & \textbf{Verdict} & \textbf{Reads as} \\
\midrule
$-3$ & acceptable & \$1 per \$1{,}000 \\
$-5$ & good & \$1 per \$100{,}000 \\
$-7$ & excellent & \$1 per \$10 million \\
\bottomrule
\end{tabular}
\end{center}
\medskip
\begin{takeawaybox}[title=\textbf{Reporting standard for this course}]\footnotesize
State the grid, the tolerance, the dtype, and the \textbf{maximum} residual on points that were \emph{not} used to solve --- not the mean, and never a residual evaluated on the solution grid itself.
\end{takeawaybox}
\end{column}
\end{columns}
\vspace{1mm}
\src{Thresholds from Judd (1992) as reported in QM Lec.~7; the separation-of-grids requirement from QM Lec.~8, ``The Evaluation Principle''.}
\end{frame}

%=============================================================================
\begin{frame}{Why a Residual You Can Compute Bounds an Error You Cannot}
\begin{columns}[T]
\begin{column}{0.53\textwidth}
\small
You never observe $\lVert g-\hat g\rVert$: knowing it would mean knowing $g$. Two theorems say the things you \emph{can} compute stand in for it.
\medskip

\textbf{Santos \& Vigo-Aguiar (1998).} With grid spacing $h$,
\[
  \lVert v-v^{h}\rVert\le \tfrac{M}{1-\beta}h^{2},
  \qquad
  \lVert g-g^{h}\rVert\le \big(\tfrac{M}{1-\beta}\big)^{1/2}h .
\]
The value converges \textbf{one order faster} than the policy --- because at the optimum the first-order condition holds, so choosing wrongly by $\epsilon$ costs value only $\epsilon^{2}$.
\medskip

\textbf{Santos (2000).} The Euler residual bounds the policy error:
$\lVert g-\hat g\rVert\le C\,\lVert\mathcal{E}\rVert$, with $C$ set by the \emph{curvature} of $u$.
\end{column}
\begin{column}{0.44\textwidth}
\begin{cautionbox}[title=\textbf{The consequence people miss}]\footnotesize
A beautifully converged value function is \textbf{not} evidence of a good policy. $\lVert v_{n+1}-v_n\rVert<10^{-9}$ is compatible with a policy wrong in the second digit --- and in 3.A2 it will be.
\end{cautionbox}
\medskip
\begin{takeawaybox}[title=\textbf{What this licenses}]\footnotesize
The residual is \emph{computable}; the policy error is not. Santos is the bridge that lets you report the first as evidence about the second --- without it, a small residual is just a small number.\\[1mm]
$C$ grows as the problem flattens, so in a nearly linear model a small residual certifies \textbf{less}: an order of magnitude, not an error bar.
\end{takeawaybox}
\end{column}
\end{columns}
\end{frame}

%=============================================================================
\begin{frame}{The Benchmark We Never Let Go Of}
\begin{columns}[T]
\begin{column}{0.52\textwidth}
\small
Set $u(c)=\ln c$, $y=zk^{\alpha}$ and $\delta=1$. Then the model has a closed form:
\[
  g(k,z)=\alpha\beta z k^{\alpha}, \qquad c(k,z)=(1-\alpha\beta)zk^{\alpha},
\]
\[
  v(k,z)=A(z)+\frac{\alpha}{1-\alpha\beta}\ln k .
\]
Full depreciation is economically silly at an annual frequency. That is not the point. The point is that we know the answer, \textbf{exactly}, at every state.
\end{column}
\begin{column}{0.45\textwidth}
\begin{takeawaybox}[title=\textbf{The rule from 1.B1}]\small
Every numerical method in this course is tested here first. \emph{If a method cannot recover $k'=\alpha\beta z k^{\alpha}$, do not trust it on anything harder.}
\end{takeawaybox}
\medskip
\begin{examplebox}\footnotesize
It is also the first place your AI assistant gets graded. A plausible-looking solver that misses a known answer by 1\% is not a small problem --- it is a wrong program.
\end{examplebox}
\end{column}
\end{columns}
\vspace{1mm}
\src{Brock \& Mirman (1972). Calibration used throughout Session 3: $\alpha=0.36$, $\beta=0.95$, $\delta=1$.}
\end{frame}

%=============================================================================
\begin{frame}{What Computing Teaches That the Theorem Does Not}
\begin{columns}[T]
\begin{column}{0.49\textwidth}
\small
The contraction theorem says $\lVert v_n-v^*\rVert\le\beta^n\lVert v_0-v^*\rVert$ and stops there. Running the algorithm tells you things the theorem is silent about:
\medskip
\begin{itemize}\footnotesize
  \item \textbf{The policy converges long before the value does.} On our grid the policy stops changing at iteration \textbf{11}; the value needs \textbf{272}.
  \item \textbf{The error has a shape, and iterating does not remove it.} The policy error is a \emph{sawtooth} of amplitude about half a grid spacing, because the policy can only land \emph{on} a grid point.
  \item \textbf{The stopping rule is not the error.} Successive iterates can differ by $10^{-6}$ while the answer is still $19\times$ further away.
\end{itemize}
\end{column}
\begin{column}{0.48\textwidth}
\begin{conceptbox}[title=\textbf{Why this matters}]\small
Every one of these is a decision you would get wrong by reasoning from the theorem alone: how long to iterate, where to put grid points, what to report.
\end{conceptbox}
\medskip
\begin{examplebox}\footnotesize
All three are measured in the next deck, on the model above, with the closed form as the judge.
\end{examplebox}
\end{column}
\end{columns}
\end{frame}

%=============================================================================
\begin{frame}{Takeaways}
\begin{takeawaybox}
\begin{itemize}\small
  \item Existence theorems are not answers. The unknown is a \textbf{function}, and solving means choosing a finite description of it plus a finite procedure that kills the residual.
  \item Both roads out of Session 1 --- Bellman and Euler --- are \textbf{constructive}. The proof technique \emph{is} the algorithm, and it carries the error bound with it.
  \item There are exactly \textbf{six decisions} to make. Sessions 4 and 5 are about making each of them better than we will today.
  \item Error comes from the grid, the tolerance, and the \textbf{arithmetic}. They add, so the smallest one is wasted effort. Use float64 for anything with an accounting identity in it.
  \item A computed number is a \textbf{claim}. Report the maximum unit-free residual on points you did not solve on, next to a closed form whenever one exists.
\end{itemize}
\end{takeawaybox}
\end{frame}

%=============================================================================
\begin{frame}{Bridge: Make Every Choice Badly, On Purpose}
\begin{columns}[T]
\begin{column}{0.53\textwidth}
\small
We now have the list of six decisions. The fastest way to understand it is to make all six in the crudest way available, run the result, and see what it costs.
\medskip
\begin{itemize}\footnotesize
  \item Values on a uniform grid, linear interpolation, the maximum found by searching the grid, a Markov chain handed to us, and a naive stopping rule.
  \item Roughly \textbf{25 lines of Python}, and we know the exact answer it has to reproduce.
\end{itemize}
\medskip
Then we count: how long it took, how wrong it is, and \emph{which of the six crude choices} is responsible for each part of the damage.
\end{column}
\begin{column}{0.44\textwidth}
\begin{conceptbox}[title=\textbf{Next (3.A2)}]\small
Value function iteration on a grid: the first complete crossing from the equation to a number.\\[1mm]
It works --- and then it tells us exactly which numerical techniques we are missing, which is the agenda for the rest of the month.
\end{conceptbox}
\end{column}
\end{columns}
\end{frame}

\end{document}
